\documentclass[10pt,a4paper]{amsart}
\usepackage{amsmath,amssymb,amsthm,graphicx}
% Retained minimal preamble for the Stacks Project chapter varieties.tex.
% Verbatim excerpt of the chapter preamble: Xy-pic setup (native preamble lines 14-18), the native theorem
% declarations with the mathematical macros (native preamble lines 154-217), and the project own no-op
% definitions of the reference/slogan/history display environments (native wrapper.tex lines 1-3) that the
% counted statement uses for its slogan. Copied mechanically; nothing added, removed or redefined.
% For commutative diagrams we use Xy-pic
\usepackage[all]{xy}

% We use 2cell for 2-commutative diagrams.
\xyoption{2cell}
% Theorem environments.
%
\theoremstyle{plain}
\newtheorem{theorem}[subsection]{Theorem}
\newtheorem{proposition}[subsection]{Proposition}
\newtheorem{lemma}[subsection]{Lemma}

\theoremstyle{definition}
\newtheorem{definition}[subsection]{Definition}
\newtheorem{example}[subsection]{Example}
\newtheorem{exercise}[subsection]{Exercise}
\newtheorem{situation}[subsection]{Situation}

\theoremstyle{remark}
\newtheorem{remark}[subsection]{Remark}
\newtheorem{remarks}[subsection]{Remarks}

\numberwithin{equation}{subsection}

% Macros
%
\def\lim{\mathop{\mathrm{lim}}\nolimits}
\def\colim{\mathop{\mathrm{colim}}\nolimits}
\def\Spec{\mathop{\mathrm{Spec}}}
\def\Hom{\mathop{\mathrm{Hom}}\nolimits}
\def\Ext{\mathop{\mathrm{Ext}}\nolimits}
\def\SheafHom{\mathop{\mathcal{H}\!\mathit{om}}\nolimits}
\def\SheafExt{\mathop{\mathcal{E}\!\mathit{xt}}\nolimits}
\def\Sch{\mathit{Sch}}
\def\Mor{\mathop{\mathrm{Mor}}\nolimits}
\def\Ob{\mathop{\mathrm{Ob}}\nolimits}
\def\Sh{\mathop{\mathit{Sh}}\nolimits}
\def\NL{\mathop{N\!L}\nolimits}
\def\CH{\mathop{\mathrm{CH}}\nolimits}
\def\proetale{{pro\text{-}\acute{e}tale}}
\def\etale{{\acute{e}tale}}
\def\QCoh{\mathit{QCoh}}
\def\Ker{\mathop{\mathrm{Ker}}}
\def\Im{\mathop{\mathrm{Im}}}
\def\Coker{\mathop{\mathrm{Coker}}}
\def\Coim{\mathop{\mathrm{Coim}}}

% Boxtimes
%
\DeclareMathSymbol{\boxtimes}{\mathbin}{AMSa}{"02}

%
% Macros for moduli stacks/spaces
%
\def\QCohstack{\mathcal{QC}\!\mathit{oh}}
\def\Cohstack{\mathcal{C}\!\mathit{oh}}
\def\Spacesstack{\mathcal{S}\!\mathit{paces}}
\def\Quotfunctor{\mathrm{Quot}}
\def\Hilbfunctor{\mathrm{Hilb}}
\def\Curvesstack{\mathcal{C}\!\mathit{urves}}
\def\Polarizedstack{\mathcal{P}\!\mathit{olarized}}
\def\Complexesstack{\mathcal{C}\!\mathit{omplexes}}
% \Pic is the operator that assigns to X its picard group, usage \Pic(X)
% \Picardstack_{X/B} denotes the Picard stack of X over B
% \Picardfunctor_{X/B} denotes the Picard functor of X over B
\def\Pic{\mathop{\mathrm{Pic}}\nolimits}
\def\Picardstack{\mathcal{P}\!\mathit{ic}}
\def\Picardfunctor{\mathrm{Pic}}
\def\Deformationcategory{\mathcal{D}\!\mathit{ef}}
\newenvironment{reference}{}{}
\newenvironment{slogan}{}{}
\newenvironment{history}{}{}
\usepackage{cleveref}
\crefname{originaldefinition}{original definition}{original definitions}
\Crefname{originaldefinition}{Original definition}{Original definitions}
\makeatletter
\newlabel{divisors-definition-zero-scheme-s}{{02OQ}{source}{}{sourcepointer-1}{}}
\newlabel{divisors-definition-zero-scheme-s@cref}{{[originaldefinition][0][]02OQ}{[1][0][]source}{}{}{}}
\makeatother
\crefname{originallemma}{original lemma}{original lemmas}
\Crefname{originallemma}{Original lemma}{Original lemmas}
\makeatletter
\newlabel{lemma-algebraic-scheme-dim-0}{{06LH}{source}{}{sourcepointer-2}{}}
\newlabel{lemma-algebraic-scheme-dim-0@cref}{{[originallemma][0][]06LH}{[1][0][]source}{}{}{}}
\makeatother
\title{An invertible sheaf on an algebraic scheme over an infinite field is generated by dim(X) + 1 sections}
\author{The Stacks Project Authors}
\date{stacks-project:varieties.tex:local-snapshot-20261011; GNU FDL 1.2 or later}
\begin{document}
\maketitle
\noindent\textit{Editorial note.} All delivered mathematical text is the author's own, copied from the pinned Stacks Project chapter file varieties.tex; no statement or proof step was rewritten, reordered or removed.\par\medskip
\noindent\textit{Editorial note.} The proof is complete as written and ends with the generation statement for W plus the single vector; no step is referred to the reader.\par\medskip
\noindent\textit{Editorial note.} The counted spans carry 2 navigation stubs and no bibliography citation.\par\medskip
\noindent\textit{Editorial note.} The preamble is the chapter's retained minimal preamble (sources/stacks-varieties/retained-preamble.tex lines 1-77), a verbatim excerpt of the native preamble.\par\medskip
\noindent\textit{Editorial note.} The label lemma-generated-by-dim-plus-1-sections is defined only in varieties.tex, so no cross-chapter restatement had to be weighed for this card.\par\medskip
\noindent\textit{Allowed standard prerequisites (editorial).}\begin{itemize}
\item Zero scheme of a section of an invertible sheaf. Divisors, Definition divisors-definition-zero-scheme-s, invoked at proof lines 4289-4291. Supplies the closed subscheme Z whose dimension is compared with dim(X).
\item Structure of an algebraic scheme of dimension zero. Same chapter, Lemma lemma-algebraic-scheme-dim-0, invoked at proof lines 4278-4280. Supplies that a zero dimensional algebraic scheme has finitely many points, making the base case a finite avoidance problem.
\item A finite union of proper hyperplanes in a vector space over an infinite field is proper. Standard linear algebra, used at proof lines 4280-4282 and 4287-4289. Supplies the existence of a vector on which none of finitely many functionals vanishes.
\end{itemize}\medskip
\ifdefined\hypertarget\hypertarget{sourcepointer-1}{}\fi
\noindent\textit{Source reference.} Original definition 02OQ: sources/stacks-chow/tags.txt, source lines 3525--3525. Pointer outside the counted span, invoked once for context; the load-bearing argument of the counted statement is the author's proof reproduced in full in the proof span.\par\medskip
\ifdefined\hypertarget\hypertarget{sourcepointer-2}{}\fi
\noindent\textit{Source reference.} Original lemma 06LH: sources/stacks-varieties/pinned-original-source.tex, source lines 3736--3736. Pointer outside the counted span, invoked once for context; the load-bearing argument of the counted statement is the author's proof reproduced in full in the proof span.\par\medskip
\begin{lemma}
\label{lemma-generated-by-dim-plus-1-sections}
Let $k$ be an infinite field. Let $X$ be an algebraic $k$-scheme.
Let $\mathcal{L}$ be an invertible $\mathcal{O}_X$-module.
Let $V \to \Gamma(X, \mathcal{L})$ be a linear map of $k$-vector spaces
whose image generates $\mathcal{L}$. Then there exists a subspace
$W \subset V$ with $\dim_k(W) \leq \dim(X) + 1$ which generates $\mathcal{L}$.
\end{lemma}

\begin{proof}
Throughout the proof we will use that for every $x \in X$
the linear map
$$
\psi_x : V \to \Gamma(X, \mathcal{L}) \to \mathcal{L}_x \to
\mathcal{L}_x \otimes_{\mathcal{O}_{X, x}} \kappa(x)
$$
is nonzero. The proof is by induction on $\dim(X)$.

\medskip\noindent
The base case is $\dim(X) = 0$. In this case $X$ has finitely many points
$X = \{x_1, \ldots, x_n\}$ (see for example
Lemma \ref{lemma-algebraic-scheme-dim-0}). Since $k$ is infinite
there exists a vector $v \in V$ such that $\psi_{x_i}(v) \not = 0$
for all $i$. Then $W = k\cdot v$ does the job.

\medskip\noindent
Assume $\dim(X) > 0$. Let $X_i \subset X$ be the irreducible components
of dimension equal to $\dim(X)$. Since $X$ is Noetherian there are only
finitely many of these. For each $i$ pick a point $x_i \in X_i$.
As above choose $v \in V$ such that $\psi_{x_i}(v) \not = 0$
for all $i$. Let $Z \subset X$ be the zero scheme of the image
of $v$ in $\Gamma(X, \mathcal{L})$, see
Divisors, Definition \ref{divisors-definition-zero-scheme-s}.
By construction $\dim(Z) < \dim(X)$. By induction we can find
$W \subset V$ with $\dim(W) \leq \dim(X)$ such that $W$ generates
$\mathcal{L}|_Z$. Then $W + k\cdot v$ generates $\mathcal{L}$.
\end{proof}

\end{document}
